\section*{Discussion}

Transient unilateral application of topical lidocaine to oral mucosa and
periodontal tissues increased masseter EMG amplitude and chewing regularity
during the first minute post-application while leaving chewing frequency and
rhythm entirely unchanged.
The amplitude difference decayed across subsequent bouts in a pattern
consistent with lidocaine pharmacokinetics, was present in 76\% of
participants, was independent of the counterbalancing order, and was
corroborated at the single-cycle level by a broadband time-frequency cluster.
The most parsimonious interpretation is sensorimotor gain compensation:
when oral mucosal and superficial periodontal afferents are transiently
reduced by sodium-channel blockade, motor recruitment of the jaw-closing
muscles increases, most plausibly via reduced inhibitory afference raising
the set-point of motor neuron recruitment.
Passive disinhibition of jaw-closing motor neurons and active upregulation
of efferent drive cannot be distinguished from surface EMG data alone;
this mechanistic ambiguity is acknowledged as a limitation (see below).
The pharmacokinetic time course -- lidocaine spray reaches peak mucosal
analgesia within approximately 2--5~min and loses efficacy within
15--20~min~\cite{vonGunten2019} -- and the null carryover and period effects
collectively support the attribution to drug-mediated sensory reduction
rather than to practice, expectation, or non-specific session effects
(further discussed in Limitations).

The amplitude-timing dissociation is the mechanistic centre of these results.
Within the CPG framework of masticatory control~\cite{Lund1991,Grillner2006},
the brainstem generates the fundamental jaw-movement rhythm while descending
and afferent signals modulate the amplitude of individual bursts to adjust
bite force to food properties~\cite{Turker2002,Trulsson2006}.
Periodontal afferents provide graded force feedback that calibrates the level
of motor neuron recruitment in the masseter on each closing
stroke~\cite{Lobbezoo1997}.
When this feedback is transiently reduced, the jaw-closing motor neuron pool
may be disinhibited, recruiting to a higher level.
That the chewing frequency channel remained entirely unaffected -- across four
bouts spanning 17~minutes, in both frequency and rate -- supports the
hierarchical dissociation between central timing and peripheral gain modulation,
and distinguishes the observed effect from a non-specific arousal or motor
excitability change, which would alter timing as well as amplitude.
It should be noted that, unlike spinal locomotor CPGs, the human masticatory
CPG is substantially shaped by afferent feedback~\cite{Lund1991,Trulsson2006}; the
preserved rhythm in the present study may reflect CPG robustness to
somatosensory deprivation, or rapid recalibration of the rhythm generator to
the altered sensory context -- a distinction that would require simultaneous
jaw-movement kinematics to resolve.

A bilateral exploratory analysis ($N$\,=\,31 with valid L\,+\,R signals)
examined whether the gain compensation was global or side-specific.
The amplitude increase at bout~1 was numerically larger in the
\emph{right} (chewing-side) masseter --- consistent with the primary
finding in the full sample --- while the \emph{left} (anesthetized-side gum)
masseter showed no detectable anesthesia effect across any bout
(Supplementary Fig.~\ref{fig:supp_bilateral}).
Neither side-specific comparison reached significance individually in
this subsample, so this pattern should be interpreted as directionally
consistent rather than conclusive.
Bilateral temporal synchrony --- the envelope correlation $r(L,R)$ ---
was preserved ($r$\,=\,0.81 vs.\ 0.84; all $p_{\text{FDR}}$\,>\,0.59;
condition\,$\times$\,bout interaction $p$\,=\,0.694), indicating that
the bilateral temporal structure of the masticatory pattern was not
reorganised under unilateral sensory reduction.
This non-significant directional trend -- if confirmed with sufficient
power -- would suggest a lateralised rather than system-wide gain change.
The neural mechanism by which unilateral block of the non-chewing-side mucosa
would preferentially affect the chewing-side masseter is not established;
crossed trigeminal projections within the masticatory nuclear complex may
contribute, but this question warrants dedicated investigation in a study
with bilateral monitoring and adequate power for side-specific contrasts.

These findings directly resolve an ambiguity in the parent study~\cite{espinoza2025},
which reported no behavioral or EEG anesthesia effect and attributed this to
a central cognitive mechanism for chewing-related theta entrainment that
operates independently of somatosensory feedback.
A potential criticism of that conclusion is that the topical anesthesia
produced no detectable peripheral effect, making the anesthesia comparison
uninformative.
The present results directly refute this concern: lidocaine altered masticatory
motor output at the level of the masseter.
The downstream null for cognition and EEG is then explained not by the absence
of a manipulation effect, but by the fact that the motor system compensated
sufficiently to preserve the temporal structure of the masticatory rhythm.
If the CPG output is preserved -- same frequency, same regularity, the same
rhythmic entrainment signal reaching the cortex -- theta entrainment and
the associated cognitive benefit follow unchanged.
The anesthesia null in the parent study is thus more accurately characterized
as motor homeostasis than as evidence of a mechanism fully impervious to
peripheral input: the masticatory CPG is robust to somatosensory deprivation
precisely because gain compensation restores the functional output.

The pharmacological time course aligns with the lidocaine PK profile.
The observed decay from a median amplitude difference of +30.5~\textmu{}V
at bout~1 ($\approx$3~min post-application) to $-$4.4~\textmu{}V at bout~3
($\approx$12~min) is consistent with the known mucosal analgesia window of
10--20~minutes~\cite{vonGunten2019}.
The exponential fit ($\tau$\,$\approx$\,2.9~min) reported in the Results is
illustrative only; with four group-level time points no reliable pharmacokinetic
parameter estimation is possible.
The robustness analysis (Supplementary Fig.~\ref{fig:supp_robust}) confirms
that the bout-1 effect size is stable across $N$\,=\,30--34 ($d_z$\,=\,0.65--0.69),
and the crossover confound analysis (Supplementary Fig.~\ref{fig:supp_carryover})
shows no sequence interaction, no period effect, and no residual carryover
from the anesthesia block to the placebo block.

Several limitations qualify these conclusions.
\textit{Motor drive index}: Masseter EMG amplitude is an indirect index of
motor drive; without simultaneous bite force measurement, it is not possible
to confirm whether the compensation was sufficient to maintain actual occlusal
force constant.
\textit{Incomplete sensory block}: The anesthesia was applied unilaterally and
topically; deep periodontal ligament mechanoreceptors may not have been fully
anaesthetised by spray application alone (see Introduction), leaving contralateral
receptors and deep masseter muscle spindles unaffected.
A complete bilateral infiltration block would be expected to produce a larger
and more sustained disruption.
\textit{Participant condition awareness}: Topical lidocaine produces a distinctive
oral numbness that participants almost certainly recognised as the active
condition; full blinding was not achievable with this design.
Although the pharmacokinetic washout pattern, the absence of carryover, and
the stability of the placebo condition across bouts collectively argue against
an expectation-driven account, a prospective study incorporating
belief-of-condition ratings or subjective sensory intensity scales would
resolve this residual confound.
\textit{Alternative arousal account}: A competing explanation for the
bout-1 amplitude increase is non-specific sympathetic arousal elicited by
the novel sensation of oral anaesthesia, which would also elevate masseter
tone and decay as habituation occurs.
The absence of significant timing changes (sympathetic arousal would be
expected to increase frequency as well as amplitude) and the PK-matched decay
provide indirect evidence against this account, but it cannot be definitively
excluded without measurement of autonomic state.
\textit{Individual variability}: The compensation was not universal -- 24\% of
participants showed the opposite pattern at bout~1 -- indicating individual
variation that future work should relate to periodontal sensitivity thresholds
or central plasticity.
\textit{EEG reanalysis pending}: Because the EEG time-frequency files from the
parent study have the anesthesia factor collapsed, a direct test of the proposed
link between motor compensation magnitude and the preservation of theta
entrainment awaits a dedicated reanalysis of the EEG data partitioned by
anesthesia condition.

